\chapter{控制流}
\label{ch:05}

程序光会「从上往下做」是干不了活的。它得会拐弯（条件）和转圈（循环）。

第 2 章的猜数字已经把这两种都用上了，这一章把它们补全：ChaLi 一共有三种循环：\texttt{while}、\texttt{loop} 和 \texttt{for-in}，还有一套对「同一个变量在一个表达式里被读写两次」管得很严的规则。

\section{条件：\texttt{if} / \texttt{elif} / \texttt{else}}

\par
\medskip
\begin{lstlisting}[style=chali]
String classify(int64 value)
{
    if (value < 0)
    {
        return "负数";
    }
    elif (value == 0)
    {
        return "零";
    }
    else
    {
        return "正数";
    }
}

void main()
{
    println("{} {} {}", classify(-3), classify(0), classify(8));

}
\end{lstlisting}
\par
\medskip

\par
\medskip
\begin{lstlisting}[style=out]
负数 零 正数
\end{lstlisting}
\par
\medskip

形态就是这些：

\par
\medskip
\begin{lstlisting}[style=out]
if (条件) { ... }
elif (条件) { ... }
else { ... }
\end{lstlisting}
\par
\medskip

三条要注意的：

1. 条件必须真的能当条件。

\par
\medskip
\begin{lstlisting}[style=chali]
void main()
{
    if (1)
    {
        println("x");
    }

}
\end{lstlisting}
\par
\medskip

\par
\medskip
\begin{lstlisting}[style=out]
chali: if 条件必须是 bool，现在给的是 int64（数值条件请显式写成 `(bool)值` 或比较表达式）
\end{lstlisting}
\par
\medskip

没有「非零即真」。想判断「\texttt{count} 是不是 0」，写 \texttt{if (\allowbreak{}count !=\allowbreak{} 0)\allowbreak{}}——把你想说的说出来。

2. 拼写是 \texttt{elif}，不是 \texttt{else if}，也不是 \texttt{elsif}。少一个字母都编译不过。

3. 花括号不能省。 只写一句 \texttt{if (\allowbreak{}x)\allowbreak{} y =\allowbreak{} 1;} 是编译错误。这条和 C 不一样，但换来的是「这个 \texttt{y =\allowbreak{} 1} 到底属于哪个 if」永远不会有歧义。

\begin{tip}{提示}
ChaLi 没有三目运算符。 \texttt{cond ? a :\allowbreak{} b} 是语法错误（第 3 章提过）。简单替换就用 \texttt{if /\allowbreak{} else} 赋值；分支多了，第 10 章的 \texttt{match} 更好读。可以使用 \texttt{if (\allowbreak{}条件)\allowbreak{} \{ 值一 \} else \{ 值二 \}} 条件表达式；只求值命中的分支。
\end{tip}

\subsection*{条件也能产生值}

\texttt{if} 可以出现在变量定义、赋值、函数实参和其他表达式里。只执行选中的分支，条件和分支都按短路规则求值。

\par
\medskip
\begin{lstlisting}[style=chali]
void main()
{
    int32 p = 5;
    bool odd = if (p % 2 == 1)
    {
        true
    }
    else
    {
        false
    };
    int32 category = if (p % 3 == 1)
    {
        1
    }
    elif (p % 3 == 2)
    {
        2
    }
    else
    {
        3
    };
    odd = if (p % 2 == 0)
    {
        false
    }
    else
    {
        bool value = true;
        yield value;
    };
    println("{} {}", odd, category);

}
\end{lstlisting}
\par
\medskip

\par
\medskip
\begin{lstlisting}[style=out]
1 2
\end{lstlisting}
\par
\medskip

只有结果表达式时写 \texttt{\{ value \}}，块内不加分号；有声明、赋值等语句时用 \texttt{yield value;} 给出块的结果。各分支结果类型一致，并且必须有 \texttt{else}。外层变量定义和赋值仍以分号结束。

条件选择使用 \texttt{if} 表达式。

\section{\texttt{while}：先判断，再执行}

\par
\medskip
\begin{lstlisting}[style=chali]
void main()
{
    int64 countdown = 3;
    while (countdown > 0)
    {
        println("{}", countdown);
        countdown = countdown - 1;
    }
    println("起飞！");

}
\end{lstlisting}
\par
\medskip

\par
\medskip
\begin{lstlisting}[style=out]
3
2
1
起飞！
\end{lstlisting}
\par
\medskip

条件写在括号里，先判后做。条件一开始就不成立，循环体一次都不执行。

\section{\texttt{for-in}：遍历区间}

计数循环使用 \texttt{std:\allowbreak{}:\allowbreak{}lang:\allowbreak{}:\allowbreak{}range(\allowbreak{}a,\allowbreak{} b)\allowbreak{}}，依次取得左闭右开区间中的值：

\par
\medskip
\begin{lstlisting}[style=chali]
void main()
{
    int64 total = 0;
    for (int64 k in std::lang::range(0, 5))
    {
        total = total + k;
    }
    println("{}", total);

}
\end{lstlisting}
\par
\medskip

\par
\medskip
\begin{lstlisting}[style=out]
10
\end{lstlisting}
\par
\medskip

输出为 \texttt{10}。需要包含上端点时使用 \texttt{std:\allowbreak{}:\allowbreak{}lang:\allowbreak{}:\allowbreak{}range\allowbreak{}Eq(\allowbreak{}a,\allowbreak{} b)\allowbreak{}}。循环变量只在循环中存在。整数和 \texttt{char} 直接支持区间遍历；其他值类型实现\texttt{Next\allowbreak{}Value<\allowbreak{}T>\allowbreak{}}，以 \texttt{Option<\allowbreak{}T>\allowbreak{} next\allowbreak{}Value(\allowbreak{}T end,\allowbreak{} bool inclusive)\allowbreak{}} 管理前进和终止。

\section{循环也能产生值}

\texttt{for-in} 和 \texttt{while} 可以作为赋值表达式，循环体用 \texttt{yield} 给出结果，\texttt{else} 在没有产生结果而退出循环时给出后备值。\texttt{continue} 继续下一轮，普通\texttt{break} 结束循环并进入 \texttt{else}。

\par
\medskip
\begin{lstlisting}[style=chali]
void main()
{
    int64 found = for (int64 value in std::lang::range(0, 10))
    {
        if (value > 5)
        {
            yield value;
        }
    }
    else
    {
        yield -1;
    };
    String label = while (false)
    {
        "运行"
    }
    else
    {
        "结束"
    };
    println("{} {}", found, label);

}
\end{lstlisting}
\par
\medskip

\par
\medskip
\begin{lstlisting}[style=out]
6 结束
\end{lstlisting}
\par
\medskip

输出为 \texttt{6 结束}。分支只有一个值时，可以省略 \texttt{yield} 和值后的分号。表达式必须有 \texttt{else}，并且每条正常结束的路径都产生兼容类型的值。返回 String、数组或容器时仍遵守值语义；不可复制的 opcode 对象需转移所有权。

\section{\texttt{loop}：会「带值跳出来」的循环}

\par
\medskip
\begin{lstlisting}[style=chali]
void main()
{
    int64 attempts = 0;
    loop
    {
        attempts = attempts + 1;
        if (attempts >= 3)
        {
            break;
        }
    }
    println("{}", attempts);

}
\end{lstlisting}
\par
\medskip

\par
\medskip
\begin{lstlisting}[style=out]
3
\end{lstlisting}
\par
\medskip

\texttt{loop \{ .\allowbreak{}.\allowbreak{}.\allowbreak{} \}} 就是「一直转」，它自己不会停——出口全靠 \texttt{break}。

它的特别之处在于：\texttt{loop} 是表达式，能算出一个结果。要交结果就用 \texttt{yield}：

\par
\medskip
\begin{lstlisting}[style=chali]
int64 firstEven = loop
{
    int64 candidate = nextNumber();
    if (candidate % 2 == 0)
    {
        yield candidate;
        // 这个 candidate 就是 firstEven
    }
};
\end{lstlisting}
\par
\medskip

\texttt{yield candidate;} 把值交给这个 \texttt{loop}，\texttt{loop} 再把它交给 \texttt{first\allowbreak{}Even}。

既然 \texttt{loop} 是表达式，分号该怎么写？看它站在哪个位置：

\begin{itemize}

\item 语句位置（单独一行，没人接它的结果）：\texttt{loop \{ .\allowbreak{}.\allowbreak{}.\allowbreak{} \}} 直接写，后面不补分号，和 \texttt{while}、\texttt{for}、\texttt{if} 一样。

\item 表达式位置（变量定义、赋值、函数实参、返回值）：\texttt{loop \{ .\allowbreak{}.\allowbreak{}.\allowbreak{} \}} 只是外层表达式的一部分，按外层语句的规矩收尾——\texttt{int64 first\allowbreak{}Even =\allowbreak{} loop \{ .\allowbreak{}.\allowbreak{}.\allowbreak{} \};} 这个分号属于那条定义语句。

\end{itemize}

\par
\medskip
\begin{lstlisting}[style=chali]
void main()
{
    int64 count = 0;
    loop
    {
        count = count + 1;
        if (count >= 2)
        {
            break;
        }
    }

    int64 doubled = loop
    {
        yield count * 2;
    };
    println("{} {}", count, doubled);

}
\end{lstlisting}
\par
\medskip

\par
\medskip
\begin{lstlisting}[style=out]
2 4
\end{lstlisting}
\par
\medskip

前一个 \texttt{loop} 是语句，\texttt{\}} 后面干净收住；后一个被 \texttt{doubled} 接走，分号写在定义语句末尾。语句位置多写一个 \texttt{;} 也能编译（那是一个结果被丢掉的表达式语句），但没必要。

\section{\texttt{for-in}：遍历一串东西}

\par
\medskip
\begin{lstlisting}[style=chali]
void main()
{
    List<int64> scores;
    scores.push(30);
    scores.push(55);
    scores.push(41);

    int64 best = 0;
    for (int64 score in scores)
    {
        if (score > best)
        {
            best = score;
        }
    }
    println("最高分 {}", best);

}
\end{lstlisting}
\par
\medskip

\par
\medskip
\begin{lstlisting}[style=out]
最高分 55
\end{lstlisting}
\par
\medskip

写法：

\par
\medskip
\begin{lstlisting}[style=out]
for ( 元素类型 名字 in 容器 ) { ... }
\end{lstlisting}
\par
\medskip

\texttt{List<\allowbreak{}int64>\allowbreak{}} 是一个「有序的整数表」，\texttt{push} 往末尾加一个。容器是第 11 章的主题，现在你把它当成「一串东西」就够了。

三个要点：

\begin{itemize}

\item 元素类型可以显式写成 \texttt{for (\allowbreak{}int64 score in scores)\allowbreak{}}，也可以由迭代源推导成 \texttt{for (\allowbreak{}score in scores)\allowbreak{}}。

\item 遍历出来的 \texttt{score} 是元素的副本（还记得第 4 章吗），改它不影响容器。想改容器里的元素得用引用（第 7 章）。

\item \texttt{for-in} 能遍历的不只是容器，凡是「能一个个吐出来的东西」都行，比如第 4 章的变长参数、第 17 章的迭代器。

\end{itemize}

显式 \texttt{for(\allowbreak{}let item in values)\allowbreak{}} 与省略类型的 \texttt{for(\allowbreak{}item in values)\allowbreak{}} 使用相同的元素类型推断；容器按元素类型，整数区间按端点类型推断。

\section{\texttt{break} 与 \texttt{continue}}

\par
\medskip
\begin{lstlisting}[style=chali]
void main()
{
    List<int64> scores;
    scores.push(12);
    scores.push(-5);
    scores.push(0);
    scores.push(99);

    for (int64 score in scores)
    {
        if (score < 0)
        {
            continue;
            // 负数不算，看下一个
        }
        if (score == 0)
        {
            break;
            // 到此为止，后面的都别看了
        }
        println("分数 {}", score);
    }

}
\end{lstlisting}
\par
\medskip

\par
\medskip
\begin{lstlisting}[style=out]
分数 12
\end{lstlisting}
\par
\medskip

（\texttt{0} 那一项触发了 \texttt{break}，所以 \texttt{99} 从没被打印。）

\texttt{continue} = 「这一轮到此为止，直接进下一轮」；\texttt{break} = 「整个循环到此为止」。

两个都只作用于最内层那一圈循环。想从两层里一次跳出来，得靠 \texttt{loop} 带值或者加个标志变量——ChaLi 没有「带标签的 break」。

\subsection*{指定退出哪一层循环}

\texttt{break;} 退出最近一层循环。\texttt{break level 2;} 退出两层，\texttt{break all;} 退出当前函数里所有包围它的循环。\texttt{level} 与 \texttt{all} 只在这个语法位置识别，其他地方仍是普通变量名。

\par
\medskip
\begin{lstlisting}[style=chali]
void main()
{
    int64 visited = 0;
    loop
    {
        loop
        {
            visited++;
            break level 2;
        }
        visited += 100;
    }

    int64 total = 0;
    List<int64> values;
    values.push(2);
    values.push(4);
    values.push(6);
    for (item in values)
    {
        total += item;
    }
    println("{} {}", visited, total);

}
\end{lstlisting}
\par
\medskip

\par
\medskip
\begin{lstlisting}[style=out]
1 12
\end{lstlisting}
\par
\medskip

层数只计算 \texttt{while}、\texttt{for-in} 和 \texttt{loop}，普通块及 \texttt{if}、\texttt{match} 不增加层数。数字参数必须是正整数字面量，不能用变量代替层数。\texttt{all} 只在 \texttt{break} 后的层级参数位置表示「全部循环」，普通变量仍可命名为 \texttt{all}。

层级退出不能携带结果，也不能跳出另一个函数的循环；离开路径上的清理照常执行。带值的 \texttt{yield value;} 则交给最近接收结果的值表达式，无值循环不会把这个值截走。返回表达式中禁止 \texttt{break level} / \texttt{break all}，详见第 4 章。

\subsection*{带值退出与层级写法小结}

\begin{center}
\begin{tabular}{ll}
\toprule
写法 & 含义 \\
\midrule
\texttt{yield 值;} & 从最近的值表达式（loop / match / 块）交出结果 \\
\texttt{break level N;} & 跳出 N 层循环，N 必须是正整数字面量 \\
\texttt{break all;} & 跳出当前函数内所有包围它的循环 \\
\texttt{break level all;} & 同 \texttt{break all;} \\
\bottomrule
\end{tabular}
\end{center}

\par
\medskip
\begin{lstlisting}[style=chali]
void main() {
    // yield 值;（写在接收结果的值表达式里）
    int64 v = loop { yield 42; };
    assert(v == 42);
    // break level N;
    int64 n = 0;
    loop { loop { n++; break level 2; } assert(false); }
    assert(n == 1);
    // break all;
    int64 a = 0;
    loop { loop { a++; break all; } assert(false); }
    assert(a == 1);
    println("break 写法 ok");

}
\end{lstlisting}
\par
\medskip

\par
\medskip
\begin{lstlisting}[style=out]
break 写法 ok
\end{lstlisting}
\par
\medskip

\texttt{yield} 只在带值退出里是关键字：作为普通标识符仍可用（例如用户自定义的 \texttt{yield(\allowbreak{})\allowbreak{}} 不受影响），但不能写成 \texttt{yield(\allowbreak{}参数)\allowbreak{}} 这种调用形式。\texttt{break level} 后面必须跟正整数字面量层数或 \texttt{all}；\texttt{break level foo;}、\texttt{break level 0;}、\texttt{break level -1;} 都是错的。\texttt{all} 仍是普通变量名，只有在这里的层级参数位置才被识别为「全部循环」，所以你照样可以写 \texttt{int64 all =\allowbreak{} 0;}。

\section{复合赋值与自增减}

写累加的时候，\texttt{total =\allowbreak{} total + score;} 说三遍就烦了。可以缩短：

\par
\medskip
\begin{lstlisting}[style=chali]
void main()
{
    int64 value = 17;
    value += 3;
    // value = value + 3
    value -= 1;
    value *= 2;
    println("{}", value);

}
\end{lstlisting}
\par
\medskip

\par
\medskip
\begin{lstlisting}[style=out]
38
\end{lstlisting}
\par
\medskip

\texttt{+=\allowbreak{}} \texttt{-=\allowbreak{}} \texttt{*=\allowbreak{}} \texttt{/\allowbreak{}=\allowbreak{}} \texttt{\%=\allowbreak{}} \texttt{\&=\allowbreak{}} \texttt{|=\allowbreak{}} \texttt{\textasciicircum{}=\allowbreak{}} \texttt{<\allowbreak{}<\allowbreak{}=\allowbreak{}} \texttt{>\allowbreak{}>\allowbreak{}=\allowbreak{}} 全都有，含义就是「把右边算出来，跟左边做这个运算，写回左边」。

单独的自增自减也能写：

\par
\medskip
\begin{lstlisting}[style=chali]
void main()
{
    int64 i = 5;
    i++;
    println("{}", i);
    i--;
    ++i;
    println("{}", i);

}
\end{lstlisting}
\par
\medskip

\par
\medskip
\begin{lstlisting}[style=out]
6
6
\end{lstlisting}
\par
\medskip

\subsection*{但 ChaLi 对「一个表达式里读写同一个对象」管得很严}

这是 ChaLi 和 C 差异最大的地方之一。同一个完整表达式里，一个自增减的目标不能再被别的操作读写：

\par
\medskip
\begin{lstlisting}[style=chali]
void main()
{
    int64 i = 0;
    i = i++;

}
\end{lstlisting}
\par
\medskip

\par
\medskip
\begin{lstlisting}[style=out]
chali: [CHL1301] 同一表达式中，自增减目标与其他读写冲突：i
\end{lstlisting}
\par
\medskip

C 程序员会心头一紧——\texttt{i =\allowbreak{} i++} 在 C 里是未定义行为（每个编译器给的结果可能不一样）。ChaLi 的做法是：干脆不让你写。你可能觉得它啰嗦，但它拦下的是一整类「换个编译器就换结果」的玄学 bug。

那什么情况是合法的？目标互不影响就行，比如两个不同的变量：

\par
\medskip
\begin{lstlisting}[style=chali]
void main()
{
    int64 i = 0;
    int64 j = 0;
    int64 sum =(i++) +(j++);
    println("{} {} {}", sum, i, j);

}
\end{lstlisting}
\par
\medskip

\par
\medskip
\begin{lstlisting}[style=out]
0 1 1
\end{lstlisting}
\par
\medskip

还有一条顺便记住：下标自增只发生一次。\texttt{xs[\allowbreak{}i]\allowbreak{} +=\allowbreak{} 1} 里，\texttt{xs[\allowbreak{}i]\allowbreak{}} 这个「位置」只算一次，不会因为复合赋值就把 \texttt{i} 加两次。

\par
\medskip
\begin{lstlisting}[style=chali]
void main()
{
    List<int64> store;
    store.push(10);
    store.push(20);
    int64 i = 0;
    store[i] += 1;
    println("{} {}", store[0], i);

}
\end{lstlisting}
\par
\medskip

\par
\medskip
\begin{lstlisting}[style=out]
11 0
\end{lstlisting}
\par
\medskip

\section{小结}

{\small\setlength{\tabcolsep}{4pt}
\begin{longtable}{>{\raggedright\arraybackslash}p{\dimexpr 0.219\linewidth-2\tabcolsep\relax}>{\raggedright\arraybackslash}p{\dimexpr 0.209\linewidth-2\tabcolsep\relax}>{\raggedright\arraybackslash}p{\dimexpr 0.552\linewidth-2\tabcolsep\relax}}
\toprule
结构 & 什么时候用 & 备注 \\
\midrule
\endhead
\texttt{if /\allowbreak{} elif /\allowbreak{} else} & 条件分支 & 条件必须是 \texttt{bool}；花括号不能省 \\
\texttt{while (\allowbreak{}cond)\allowbreak{} \{ \}} & 先判断后执行 & 可能一次都不执行 \\
\texttt{loop \{ .\allowbreak{}.\allowbreak{}.\allowbreak{} yield v; \}} & 出口不确定、要给个结果 & 是表达式；语句位置不补分号，表达式位置按外层语句收尾 \\
\texttt{for (\allowbreak{}T x in c)\allowbreak{} \{ \}} & 遍历一串东西 & 类型可显式写出或推导，遍历出来是副本 \\
\texttt{break} / \texttt{continue} & 跳出 / 下一个 & 普通无值形式作用于最近循环；\texttt{break level N} 可指定退出层数 \\
\bottomrule
\end{longtable}
}

还有一套要记住的规则：同一个表达式里对同一个目标读写两次是被禁止的 （\texttt{CHL1301}），这是 ChaLi 用来消灭「未定义行为」的手段。

三件本章没讲的事：

\begin{itemize}

\item \texttt{for-in} 遍历出来的元素是副本，怎么改到原来的——第 7 章；

\item 分支多了 \texttt{if /\allowbreak{} elif} 会很难看，\texttt{match} 怎么救场——第 10 章；

\item \texttt{loop} 里 \texttt{yield} 出来的值，除了赋值还能怎么用——第 10 章的 \texttt{match}、第 16 章的 \texttt{?} 会陆续出现。

\end{itemize}

\section{下一步}

到这里，「语言基础」这一部分就结束了。你现在会写变量、函数、判断、循环——用这些已经能写出完整的程序。

但会写程序 \ensuremath{\neq} 会写对的程序。接下来三章是这本书最重要的一段：【第 6 章 值语义】【第 7 章 引用】【第 8 章 资源与作用域】。

读懂了这三章，你会明白 ChaLi 为什么长成现在这样——以及为什么它愿意让你多写几个字。

翻到\hyperref[ch:06]{第 6 章 值语义}。

\section{边界与验证}

break 和 continue 的目标由语法和所在循环决定。修改嵌套结构后必须重新核对目标；带值退出还需要满足所在表达式的结果类型要求。

\section{习题}

前两题用于理解和基本实现，第三题检查边界，最后一题为综合挑战。先独立完成，再核对参考解答。

\begin{enumerate}

\item 解释内层循环的 break 对外层循环有什么影响。

\item 输出 1 到 20 中所有偶数，并验证输出项数。

\item 实现两层查找，并在发现目标后结束整个查找。

\item 挑战：将查找结果写成带值的循环表达式，说明所有退出路径的类型。

\end{enumerate}

\subsection*{参考解答与验收要点}

\begin{enumerate}

\item 普通 break 退出当前对应的最近循环，外层通常继续；不能按缩进猜测为退出函数。

\item 用范围内的遍历和余数判断，预期十项；检查端点是否包含。

\item 采用本章讲过的退出方式或明确状态变量；同时测试目标不存在的路径。

\item 成功和失败路径的值必须与表达式类型相容；不要遗漏无匹配时的终止路径。

\end{enumerate}

\section{番外：桃花源里的 break}

改编自陶渊明《桃花源记》，以下为编程主题的独立改写。

渔人进了桃花源，村人请他修一段寻找出口的程序。外层逐村找路，内层逐户问人，找到写着“出口”的门便 break。

他按下运行，终于走出那一户，转眼又站到下一村门前。村人赞叹程序并未死掉。渔人却说自己要离开整片桃花源，不是离开一间屋。

老人取来纸，把两层循环画成两个院子，请渔人指出要退出哪一个。渔人这才认出，先前那声 break 只说了最近的一道门。他重新安排退出的位置，也为没找到出口的情形留了返回结果。

临走时，他在路上刻下标记。后来再来的人照着走，却仍须先看当前站在哪一层循环里。纸上的门再清楚，也不会替一个没有写出的条件自行打开。

\subsection*{技术对照}

break 和 continue 的目标由语法和所在循环决定。修改嵌套结构后必须重新核对目标；带值退出还需要满足所在表达式的结果类型要求。
